
This paper examined whether the pre-alignment confidence margin --- computable from a base model before any alignment procedure --- predicts post-alignment argmax instability. For a DPO-aligned tulu-2-7b model pair, a mixed-effects logistic regression yields $\beta _1 = -4.33$ (p $\approx 10^{-227}, \eta^2$ = 0.289), and AUROC is 0.867 (MMLU), 0.803 (TruthfulQA), and 0.909 (ARC-Challenge) with no domain adaptation between benchmarks. These results support the existence of a predictive geometric signal in the pre-alignment base model distribution for this pair. A secondary model pair (pythia-6.9b SFT) shows the directionally consistent but weaker result (AUROC = 0.609).

Alignment-induced logit perturbations are non-isotropic under both DPO (anisotropy ratio 2.90, p = 0.003 on MMLU) and SFT (4.58, p = 0.005 on MMLU), compared to an isotropic Gaussian control of 1.13. This is consistent with the geometric interpretation of the margin predictor: structured perturbations along preferred directions would disproportionately affect low-margin items near the argmax boundary.

The H-M2 directional hypothesis --- that DPO concentrates perturbation variance in low-confidence quintiles --- is not supported; the direction is reversed across all three benchmarks. DPO shows a monotone increasing quintile variance profile (Q5/Q1 ratio = 4.$79\times )$ while SFT is flat (Q5/Q1 ratio $\approx$ 1.$26\times )$. This null result is documented alongside the descriptive finding about the contrasting quintile profiles.

Important limitations remain: no PPO-aligned model pairs were available, the primary evidence for the margin-flip result comes from a single model pair, two planned mechanistic analyses (H-M3, H-M4) were not completed, and all results are restricted to 4-option MCQ benchmarks. A same-base-model DPO vs. SFT comparison and replication across additional model families are the most direct next steps.

